% !TEX program = xelatex
% ============================================================
%  خطة إعادة تصميم صفحة Hero — تجربة جوال جذابة
%  الاستخدام على Overleaf:
%    1) New Project ← Blank Project
%    2) الصق المحتوى في main.tex
%    3) Menu ← Compiler ← اختر: XeLaTeX   (مهم للعربية)
%    4) Recompile ← نزّل PDF
% ============================================================
\documentclass[11pt]{article}

\usepackage{polyglossia}
\setmainlanguage{arabic}
\setotherlanguage{english}
\newfontfamily\arabicfont[Script=Arabic,Scale=1.0]{Amiri}
\newfontfamily\arabicfontsf[Script=Arabic,Scale=1.0]{Amiri}
\setmonofont{DejaVu Sans Mono}[Scale=0.8]

\usepackage{geometry}
\geometry{a4paper, margin=2.2cm}

\usepackage{xcolor}
\usepackage{graphicx}
\usepackage{booktabs}
\usepackage{longtable}
\usepackage{array}
\usepackage{enumitem}
\usepackage{titlesec}
\usepackage{fancyhdr}
\usepackage{hyperref}
\usepackage{tcolorbox}
\tcbuselibrary{skins,breakable}
\usepackage{fontawesome5}

% ── الألوان (هوية يَسِّر) ──
\definecolor{gold}{HTML}{B8933B}
\definecolor{darkbg}{HTML}{101414}
\definecolor{okgreen}{HTML}{27AE60}
\definecolor{warnred}{HTML}{C0392B}
\definecolor{softgray}{HTML}{F4F4F4}
\definecolor{midgray}{HTML}{7A7A7A}
\definecolor{codebg}{HTML}{1A2020}

\pagestyle{fancy}
\fancyhf{}
\renewcommand{\headrulewidth}{0.4pt}
\renewcommand{\footrulewidth}{0.4pt}
\fancyhead[R]{\small\color{gold}خطة تصميم صفحة Hero — منصة يَسِّر}
\fancyhead[L]{\small\color{midgray}Mobile-First Redesign}
\fancyfoot[C]{\small\thepage}

\titleformat{\section}{\Large\bfseries\color{gold}}{\thesection}{0.6em}{}
\titleformat{\subsection}{\large\bfseries\color{darkbg}}{\thesubsection}{0.5em}{}
\titleformat{\subsubsection}{\normalsize\bfseries\color{midgray}}{\thesubsubsection}{0.4em}{}

\setlist[itemize]{rightmargin=0pt, itemsep=2pt, topsep=3pt}
\setlist[enumerate]{rightmargin=0pt, itemsep=2pt, topsep=3pt}

\hypersetup{
  colorlinks=true, linkcolor=gold, urlcolor=gold,
  pdftitle={خطة تصميم صفحة Hero - منصة يَسِّر}, pdfauthor={فريق يَسِّر}
}

\newtcolorbox{infobox}[1][]{colback=softgray,colframe=gold,boxrule=0.6pt,arc=3pt,left=8pt,right=8pt,top=6pt,bottom=6pt,breakable,#1}
\newtcolorbox{riskbox}[1][]{colback=red!4,colframe=warnred,boxrule=0.6pt,arc=3pt,left=8pt,right=8pt,top=6pt,bottom=6pt,breakable,#1}
\newtcolorbox{okbox}[1][]{colback=green!4,colframe=okgreen,boxrule=0.6pt,arc=3pt,left=8pt,right=8pt,top=6pt,bottom=6pt,breakable,#1}

\newcommand{\code}[1]{\textcolor{darkbg}{\texttt{\small #1}}}

\begin{document}

% ===================== الغلاف =====================
\begin{titlepage}
\centering
\vspace*{2.5cm}
{\Huge\bfseries\color{gold} يَسِّر}\\[0.4cm]
{\large\color{midgray} منصة تعليم الفقه الإسلامي}\\[2.5cm]

{\Huge\bfseries خطة إعادة تصميم}\\[0.5cm]
{\Huge\bfseries صفحة البداية (Hero)}\\[1cm]
{\LARGE\bfseries تجربة جوال جذّابة — Mobile-First}\\[1.2cm]

\rule{0.7\textwidth}{0.5pt}\\[0.8cm]
{\large Hero Page Redesign Roadmap}\\[0.8cm]
\rule{0.7\textwidth}{0.5pt}\\[2.5cm]

\begin{tabular}{r l}
\textbf{الهدف:}       & صفحة بداية جميلة وجذّابة خاصة للفون \\
\textbf{المبدأ:}      & هوية ← قيمة ← فعل (Mobile-First) \\
\textbf{الإطار:}      & عربي RTL · صفر سكرول · عند احترام تقليل الحركة \\
\textbf{المدة:}       & ≈ ٣.٥ أيام عمل \\
\textbf{النسخة:}      & 1.0 \\
\end{tabular}

\vfill
{\small\color{midgray} وثيقة تخطيط تصميمي — \today}
\end{titlepage}

\newpage
\tableofcontents
\newpage

% ===================== الملخص =====================
\section{الملخص التنفيذي}

\begin{infobox}
إعادة تصميم صفحة البداية (Hero) بمنهج \textbf{Mobile-First}، مع الحفاظ على روح المنصة الإسلامية المعاصرة.
المبدأ: \textbf{«هوية أولًا، ثم قيمة، ثم فعل»}. الخطوة الأولى \textbf{تبسيط الصفحة} بحذف بطاقات الكتب ووسوم المزايا
(كتل ثقيلة تشتّت على الجوال)، ثم أهم تغيير جذري هو \textbf{إظهار المحراب/الشعار على الجوال}
(كان مخفيًا بالكامل) ليمنح الصفحة بؤرة بصرية فورية، مع نظام أبعاد متجاوب بالارتفاع يمنع القص، وتسلسل حركة دخول احترافي.
\end{infobox}

\subsection*{تشخيص الجوال الحالي — نقاط الضعف}
\begin{center}
\renewcommand{\arraystretch}{1.35}
\begin{tabular}{@{}p{6cm} p{8cm}@{}}
\toprule
\textbf{الملاحظة} & \textbf{الأثر} \\
\midrule
الشعار مخفي تمامًا على الجوال & لا هوية بصرية مركزية — أكبر فرصة مهدورة \\
محتوى ثابت + \code{overflow:hidden} + \code{100dvh} & خطر قص/تراكم على الشاشات القصيرة \\
ترتيب عمودي مسطّح (نص ثم بطاقات) & لا تدرّج بصري ولا «مفاجأة» \\
دخول ثابت بلا حركة & لا إحساس «احترافي» \\
بطاقات الكتب ووسوم المزايا & ثقيلة نصًّا وتشوّش على شاشة صغيرة \\
\bottomrule
\end{tabular}
\end{center}

% ===================== المراحل =====================
\section{المراحل}

\subsection{المرحلة ٠ — حذف بطاقات الكتب ووسوم المزايا (قرار تصميمي)}
قرار صريح: تبسيط الصفحة بإزالة الكتل الثقيلة التي تسبّب تشتّتًا على الجوال، وإخلاء المساحة لهوية الصفحة وأدواتها الأساسية.

\begin{center}
\renewcommand{\arraystretch}{1.3}
\begin{tabular}{@{}p{6cm} p{8cm}@{}}
\toprule
\textbf{ما يُحذَف} & \textbf{السبب} \\
\midrule
بطاقتا «كتاب الطهارة» و«كتاب الصلاة» & ثقيلتان نصًّا، وتكرار لمنطقة «تصفّح الكتب» \\
وسم «٨٥ مسألة مُحققة» & رقم ثابت يتقادم مع تغيّر المحتوى \\
وسم «شرح صوتي متقن» & تفصيل يُدرك داخل المسألة لا في المقدمة \\
وسم «مزامنة سحابية فورية» & تفصيل تقني لا يعنيه في أول نظرة \\
\bottomrule
\end{tabular}
\end{center}

\begin{riskbox}
\textbf{الأثر على الملفات:} يجب إزالة كل الأنماط المرتبطة (\code{books-portal-grid}، \code{book-portal-card}، \code{book-card-*}، \code{features-strip}، \code{feature-chip}، \code{chip-icon}) ومن كل \code{media queries}، وخصائص المكوّن (\code{onSelectBook}، \code{booksStats})، والاستيرادات غير المستخدمة، ثم \code{typecheck}/\code{lint}/\code{knip} للتأكد.
\end{riskbox}

\begin{okbox}
\textbf{الباقي في الصفحة بعد الحذف:} الشريط العلوي · الوسام التعريفي · عنوان «يَسِّر» · الوصف · زر CTA الرئيسي · زران ثانويان (PDF والسلسلة المرئية) · محراب الشعار · بطاقة «آخر وقفة» · الحديث الشريف.
\end{okbox}

\subsection{المرحلة ١ — إعادة هيكل المحتوى (الأهم)}
الترتيب الجديد يتدرّج بصريًا: هوية ← قيمة ← فعل.

\begin{tcolorbox}[colback=softgray, colframe=gold, boxrule=0.6pt, arc=3pt, left=10pt, right=10pt]
\small\raggedleft
\begin{itemize}[label={\textcolor{gold}{\textbullet}}, leftmargin=1.4em, itemsep=1pt]
\item شريط علوي رفيع: شعار مصغّر + زر الدخول.
\item \textbf{محراب الشعار} (الشعار ≈٩٦px) بتوهّج ذهبي نبضي.
\item عنوان «يَسِّر» بخط Aref Ruqaa وتدرّج ذهبي.
\item سطر الوصف: «الفقه الميسر في رداءٍ عصري».
\item \textbf{زر CTA رئيسي} عريض: «تصفّح أبواب الفقه الآن».
\item زران ثانويان: «تحميل PDF» و«السلسلة المرئية».
\item بطاقة «متابعة آخر وقفة» (إن وُجدت).
\item شريط الحديث الشريف في الأسفل.
\end{itemize}
\end{tcolorbox}

\begin{okbox}
\textbf{التغيير الجوهري:} إظهار الشعار على الجوال داخل «محراب مصغّر» بدل إخفائه.
\end{okbox}

\subsection{المرحلة ٢ — نظام الأبعاد المتجاوب بالارتفاع}
المشكلة: \code{overflow:hidden} مع محتوى ثابت يسبب القص على الشاشات القصيرة.
الحل: نظام \code{clamp()} محسوب بالـ\code{dvh} بدل البكسل الثابت.

\begin{center}
\renewcommand{\arraystretch}{1.3}
\begin{tabular}{@{}p{5cm} p{8cm}@{}}
\toprule
\textbf{العنصر} & \textbf{القيمة المقترحة} \\
\midrule
ارتفاع الشعار & \code{clamp(72px, 13dvh, 108px)} \\
عنوان «يَسِّر» & \code{clamp(2.8rem, 9dvh, 4.2rem)} \\
كل الفجوات & \code{clamp(., .dvh, .)} \\
\bottomrule
\end{tabular}
\end{center}

\subsubsection*{سلّم أولويات الإخفاء عند ضيق الارتفاع}
\begin{enumerate}
\item أكبر من ٧٤٠: كل شيء ظاهر.
\item ٦٠٠–٧٤٠: أخفِ الوصف النصي الطويل.
\item أقل من ٦٠٠: أخفِ وسوم المزايا + صغّر الشعار.
\item أقل من ٥٢٠: أخفِ الحديث الشريف (يبقى في شاشة أخرى).
\end{enumerate}

\subsection{المرحلة ٣ — اللمسات البصرية}
\begin{center}
\renewcommand{\arraystretch}{1.3}
\begin{tabular}{@{}p{4cm} p{10cm}@{}}
\toprule
\textbf{العنصر} & \textbf{التحسين} \\
\midrule
محراب الشعار & قوس إسلامي بتدرّج ذهبي + توهّج نبضي خفيف \\
عنوان «يَسِّر» & خط Aref Ruqaa + تدرّج ذهبي على النص + ظل توهّج \\
بطاقات الكتب & حواف أعلى مرفوعة (شكل كتاب) + شريط لوني جانبي \\
زر CTA & نبض/لمعان خفيف (كنمط زر التحديث المعتمد) \\
الخلفية & mesh متدرّج + توهّجان (ذهبي/زمردي) + نقش زخرفي \\
الحديث الشريف & إطار زخرفي + ✦ على الطرفين + خط Amiri \\
\bottomrule
\end{tabular}
\end{center}

\subsection{المرحلة ٤ — حركة الدخول (Choreography)}
تسلسل متدرّج (stagger) عند التحميل — كل عنصر يدخل متأخرًا قليلًا:

\begin{center}
\renewcommand{\arraystretch}{1.25}
\begin{tabular}{@{}c p{8cm}@{}}
\toprule
\textbf{التوقيت} & \textbf{العنصر} \\
\midrule
٠ms & محراب الشعار (scale 0.8 ← 1 + fade) \\
١٥٠ms & عنوان «يَسِّر» (fade + translateY) \\
٣٠٠ms & الوصف + الوسوم \\
٤٥٠ms & الزر الرئيسي (slide من الأسفل) \\
٦٠٠ms & الأزرار الثانوية \\
٧٥٠ms & بطاقات الكتب (فارق ١٠٠ms بينهما) \\
٩٠٠ms & الحديث الشريف \\
\bottomrule
\end{tabular}
\end{center}

\begin{riskbox}
استخدام \code{framer-motion} (موجود) أو CSS keyframes، مع \textbf{احترام \code{prefers-reduced-motion}} كما في بقية المشروع.
\end{riskbox}

\subsection{المرحلة ٥ — تفاعل الجوال}
\begin{center}
\renewcommand{\arraystretch}{1.3}
\begin{tabular}{@{}p{4cm} p{10cm}@{}}
\toprule
\textbf{العنصر} & \textbf{التحسين} \\
\midrule
اللمس & \code{:active} scale 0.97 على كل بطاقة (feedback فوري) \\
Safe Area & \code{env(safe-area-inset-*)} لمنع القص تحت notch \\
الشريط العلوي & تصغير عند التمرير + \code{backdrop-filter} \\
الخطوط & لا خطوط أصغر من 14px على الجوال \\
أهداف اللمس & كل زر ≥ 44px (معيار Apple) \\
التمرير & \code{overscroll-behavior: none} لمنع الارتداد \\
\bottomrule
\end{tabular}
\end{center}

\subsection{المرحلة ٦ — التحقق}
\begin{center}
\renewcommand{\arraystretch}{1.3}
\begin{tabular}{@{}p{6cm} p{6cm}@{}}
\toprule
\textbf{الفحص} & \textbf{المعيار} \\
\midrule
Pixel 7 (412×915) & بلا سكرول، بلا قص \\
iPhone SE (375×667) & الأهم — الشاشة القصيرة \\
Galaxy S8 (360×740) & متوازن \\
سطح مكتب 1440×900 & التصميم الحالي محفوظ \\
Lighthouse & \code{CLS < 0.1}، لا layout shift \\
فحص RTL & النصوص والأيقونات متسقة \\
e2e & اختبار Playwright على Pixel 7 يمر \\
\bottomrule
\end{tabular}
\end{center}

% ===================== الجدول =====================
\section{الجدول الزمني}
\begin{center}
\renewcommand{\arraystretch}{1.3}
\begin{tabular}{@{}p{8cm} c@{}}
\toprule
\textbf{المرحلة} & \textbf{المدة} \\
\midrule
٠ — حذف الكتب والمزايا & ٠.٢٥ يوم \\
١ — هيكل المحتوى & ٠.٥ يوم \\
٢ — نظام الأبعاد & ٠.٥ يوم \\
٣ — اللمسات البصرية & ١ يوم \\
٤ — الحركة & ٠.٥ يوم \\
٥ — تفاعل الجوال & ٠.٥ يوم \\
٦ — التحقق & ٠.٥ يوم \\
\midrule
\textbf{الإجمالي} & \textbf{≈ ٣.٧٥ أيام} \\
\bottomrule
\end{tabular}
\end{center}

% ===================== Done =====================
\section{معايير الإنجاز (Definition of Done)}
\begin{okbox}
\begin{itemize}
\item[\faCheck] حُذفت بطاقتا الكتب ووسوم المزايا الثلاثة، وأنماطها من كل \code{media queries}.
\item[\faCheck] حُذفت خصائص \code{onSelectBook} و\code{booksStats} من المكوّن ومن تمريرها في \code{App.tsx}.
\item[\faCheck] \code{typecheck} و\code{lint} و\code{knip} خضراء بلا استيرادات/أنماط ميتة.
\item[\faCheck] الشعار ظاهر على الجوال (محراب مصغّر).
\item[\faCheck] صفر سكرول على Pixel 7 و iPhone SE.
\item[\faCheck] تسلسل حركة دخول سلس (مع احترام \code{reduced-motion}).
\item[\faCheck] كل زر ≥ ٤٤px، و\code{:active} feedback واضح.
\item[\faCheck] Safe-area مُطبّقة (لا قص تحت notch).
\item[\faCheck] Lighthouse \code{CLS} أقل من ٠.١.
\item[\faCheck] اختبار e2e يمر.
\end{itemize}
\end{okbox}

% ===================== ملحق =====================
\section{ملحق — الملفات المتأثرة}
\begin{tcolorbox}[colback=codebg, colframe=codebg, coltext=white, arc=3pt, left=8pt, right=8pt]
\ttfamily\small\raggedleft
src/Components/Header/\\
\hspace*{1em}HeroSection.tsx\hspace{1em}\textcolor{okgreen}{مُعدّل — حذف الكتب/المزايا + هيكل الجوال}\\
\hspace*{1em}HeroSection.module.css\hspace*{1em}\textcolor{okgreen}{مُعدّل — حذف الأنماط + الأبعاد}\\
src/App.tsx\hspace*{1em}\textcolor{okgreen}{مُعدّل — إزالة تمرير onSelectBook/booksStats}\\
\end{tcolorbox}

\vfill
\begin{center}
\rule{0.5\textwidth}{0.4pt}\\[0.4cm]
{\small\color{midgray} نهاية الوثيقة — منصة يَسِّر}
\end{center}

\end{document}
